% PDF p. 54
\AppendixSection{1A}{格拉斯曼流形的表示}
下面分步证明，格拉斯曼流形 $\Gr(k,n)$ 确实可以表示成秩为 $k$ 的对称正交投影矩阵组成的集合。

\Subheading{第 1 步\quad 格拉斯曼流形的定义}
格拉斯曼流形 $\Gr(k,n)$ 定义为 $\R^n$ 中全体 $k$ 维子空间所成的空间，
\[
  \Gr(k,n)=\{U\subseteq\R^n\mid
    U\text{ 是满足 }\dim(U)=k\text{ 的线性子空间}\}.
\]

\Subheading{第 2 步\quad 基的等价类（商空间定义）}
每个 $k$ 维子空间 $U\subseteq\R^n$ 都可通过选取一组标准正交基来表示：
\[
  U=\operatorname{span}\{u_1,\ldots,u_k\},\qquad
  u_i\in\R^n,\qquad u_i^Tu_j=\delta_{ij}.
\]
但这种表示不唯一，因为同一个子空间 $U$ 的不同标准正交基可以相差一个正交群 $O(k)$ 中的正交变换：
\[
  \begin{aligned}
    U&=\operatorname{span}\{u_1,\ldots,u_k\}
      =\operatorname{span}\{u'_1,\ldots,u'_k\},\\
    (u'_1,\ldots,u'_k)&=(u_1,\ldots,u_k)Q,\qquad Q\in O(k).
  \end{aligned}
\]
因而，每个子空间 $U$ 对应于标准正交标架的一个等价类，由此得到商结构
\[
  \Gr(k,n)=\frac{O(n)}{O(k)O(n-k)}.
\]
图~\ref{fig:1A.1} 直观展示了作为商空间 $O(3)/(O(2)\times O(1))$ 的格拉斯曼流形 $\Gr(2,3)$。每对向量（一红一蓝）都给出同一个二维子空间（平面）的标准正交基。绕 $z$ 轴旋转该基，便得到不同的标准正交标架，它们在格拉斯曼流形上全部等价。这些基都代表同一个等价类或子空间。因此，格拉斯曼流形把旋转得到的各组标准正交基等同为同一个子空间，直观地展示了商结构。

\Subheading{第 3 步\quad 正交投影算子}
可以用一组标准正交基构造正交投影算子。考虑到给定 $k$ 维子空间 $U$ 的正交投影算子 $P_U$。若 $\{u_1,\ldots,u_k\}$ 是 $U$ 的一组标准正交基，则定义
\[
  U=\operatorname{span}\{u_1,\ldots,u_k\},\qquad
  U=[u_1,\ldots,u_k]\in\R^{n\times k}.
\]
于是，到子空间 $U$ 的正交投影矩阵 $P_U$ 为
\[
  P_U=UU^T.
\]

% PDF p. 55
\begin{figure}[htbp]
  \centering
  % PDF p. 55, figure 1A.1. The nearly vertical red line has no arrowheads in the original.
\begin{tikzpicture}[x=1cm,y=1cm,>=Stealth]
  \definecolor{quotientred}{HTML}{F0141B}
  \definecolor{quotientblue}{HTML}{006AB6}
  \draw[line width=.85pt,quotientred] (-.13,2.98) -- (.13,-3.09);
  \foreach \endpoint in {(-2.85,2.17),(-3.55,1.05),(-3.61,-.18),
      (-3.41,-1.05),(-2.72,-1.91),(2.59,1.92),(3.02,-2.24)} {
    \draw[->,line width=.85pt,quotientred] (0,0) -- \endpoint;
  }
  \foreach \endpoint in {(3.16,.95),(3.00,.15),(3.51,-.60)} {
    \draw[->,line width=.85pt,quotientblue] (0,0) -- \endpoint;
  }
\end{tikzpicture}

  \FigureCaption{1A.1}{格拉斯曼流形表示为 $O(3)$ 对 $O(2)\times O(1)$ 的商。}
\end{figure}

矩阵 $P_U$ 具有下列性质。对称性为 $P_U^T=P_U$，幂等性为 $P_U^2=P_U$，正交性（秩的性质）写为 $\rank(P_U)=k$。因此，每个子空间 $U$ 唯一确定一个秩为 $k$ 的对称投影矩阵 $P_U$。

\Subheading{第 4 步\quad 一一对应}
反过来，给定一个秩为 $k$ 的对称投影矩阵 $P$，可以把它的像作为唯一的子空间 $U$ 恢复出来。给定对称的秩 $k$ 投影矩阵 $P_U$，有
\[
  U=\operatorname{range}(P_U)=\{P_Ux\mid x\in\R^n\}.
\]
因此，子空间与投影矩阵之间存在一一对应关系。

\Subheading{第 5 步\quad 得出两种定义的等价性}
我们已经给出两种表示。
\begin{itemize}
  \item \textbf{按子空间定义格拉斯曼流形}
  \[
    \Gr(k,n)=\{\R^n\text{ 中全体 }k\text{ 维子空间}\}.
  \]
  \item \textbf{按投影矩阵定义}
  \[
    \Gr(k,n)\cong
    \{P\in\R^{n\times n}\mid P^T=P,\ P^2=P,\ \rank(P)=k\}.
  \]
\end{itemize}
这两种定义完全等价。

% PDF p. 56
\BackSection{结语}
我们证明了，起初定义为基的商空间的格拉斯曼流形，等价地表示为秩为 $k$ 的对称正交投影矩阵所成的集合，
\[
  \Gr(k,n)\cong
  \{P\in\R^{n\times n}\mid P^T=P,\ P^2=P,\ \rank(P)=k\}
  \cong\frac{O(n)}{O(k)O(n-k)}.
\]
这种等价关系很重要。投影矩阵表示便于计算，尤其适用于涉及优化和嵌入的应用，而商空间表示有助于进一步理解其几何和结构。

\BackSection{习题}
\begin{enumerate}
  \item \textbf{习题 1.1}\par 给出 $\R P^2$ 的数学定义。
  \item \textbf{习题 1.2}\par 定义乘积流形 $S^1\times S^2$。
  \item \textbf{习题 1.3}\par 在圆周 $S^1$ 上定义一张坐标卡。
  \item \textbf{习题 1.4}\par 给出流形上非光滑函数的一个例子。
  \item \textbf{习题 1.5}\par 设 $M(m\times n,\R)$ 表示全体 $m\times n$ 实矩阵组成的集合。证明 $M(m\times n,\R)$ 是一个流形。
  \item \textbf{习题 1.6}\par 证明全体 $n$ 维正交矩阵组成的集合是一个光滑流形。
  \item \textbf{习题 1.7}\par 证明 $L_1=\{(x,y)\mid x^2+y^2=1\}$ 是一个光滑流形。
  \item \textbf{习题 1.8}\par 证明 $\Gr(1,2)\cong S^1/\{\pm1\}$。
  \item \textbf{习题 1.9}\par 考虑 DNA 序列 ATTCGAGGAA，构造它的 CGR 图。假定网格为 $2\times2$，计算其 FCGR。
  \item \textbf{习题 1.10}\par 考虑 DNA 序列 ATTCGAGGAA，用扩散把确定性的 CGR 点转化成含噪分布。
\end{enumerate}

% PDF p. 57
\needspace{13\baselineskip}
\BackSection{参考文献}
\begin{bookreferences}
  \item Barnsley, M. F. (2012). \textit{Fractals Everywhere}. New Edition. Dover Publications.
  \item Bendokat, T., Zimmermann, R. and Absil, P. A. (2024). A Grassmann manifold handbook: basic geometry and computational aspects. \textit{Advances in Computational Mathematics}. 50: 6.
  \item Ho, J., Jain, A. and Abbeel, P. (2020). Denoising diffusion probabilistic models. In Larochelle, H., Ranzato, M., Hadsell, R., Balcan, M. and Lin, H., editors, \textit{Advances in Neural Information Processing Systems 33}: 6840–6851, PMLR. Annual Conference on Neural Information Processing Systems 2020 (NeurIPS 2020).
  \item Nichol, A. Q. and Dhariwal, P. (2021). Improved denoising diffusion probabilistic models. In Meila, M. and Zhang, T., editors, \textit{Proceedings of the 38th International Conference on Machine Learning, Volume 139 of Proceedings of Machine Learning Research}: 8162–8171. PMLR.
  \item Strümke, I. and Langseth, H. (2023). Lecture notes in probabilistic diffusion models. arXiv:2312.10393v1.
\end{bookreferences}
